\documentclass[a4paper,11pt]{article}
\usepackage[utf8]{inputenc}
\usepackage[T1]{fontenc}
\usepackage[spanish]{babel}
\usepackage{geometry}
\usepackage{amsmath}
\usepackage{amssymb}
\usepackage{listings}
\usepackage{xcolor}

% Configuración de márgenes
\geometry{left=2.5cm, right=2.5cm, top=2.5cm, bottom=2.5cm}

% Configuración para código R
\definecolor{codegray}{rgb}{0.5,0.5,0.5}
\definecolor{codegreen}{rgb}{0,0.6,0}
\definecolor{backcolour}{rgb}{0.95,0.95,0.92}

\lstset{
	language=R,
	backgroundcolor=\color{backcolour},
	commentstyle=\color{codegreen},
	keywordstyle=\color{blue},
	numberstyle=\tiny\color{codegray},
	stringstyle=\color{purple},
	basicstyle=\ttfamily\small,
	breakatwhitespace=false,
	breaklines=true,
	captionpos=b,
	keepspaces=true,
	numbers=left,
	numbersep=5pt,
	showspaces=false,
	showstringspaces=false,
	showtabs=false,
	tabsize=2
}

\title{\textbf{Documentación Técnica: Generador de Cuadrados Medios}}
\author{Implementación en R}
\date{\today}

\begin{document}
	
	\maketitle
	
	\section{Descripción General}
	La función \texttt{cuadrados\_medios\_refactored} implementa el algoritmo de \textit{Middle-square method} (Cuadrados Medios) propuesto por John von Neumann para la generación de números pseudoaleatorios uniformes $U \sim [0, 1)$.
	
	Esta implementación es una refactorización optimizada para el entorno R, mejorando la legibilidad del código y manteniendo la equivalencia matemática exacta con versiones anteriores (librería \texttt{simres}).
	
	\section{Uso}
	\begin{lstlisting}
		numeros <- cuadrados_medios_refactored(n, seed, k = 4, save_seed = TRUE)
	\end{lstlisting}
	
	\section{Argumentos}
	\begin{description}
		\item[n (entero)] Cantidad de números pseudoaleatorios a generar. Debe ser un entero positivo ($n > 0$).
		
		\item[seed (numérico)] Semilla inicial. Si no se especifica, se utiliza \texttt{as.numeric(Sys.time())}. La función normaliza la semilla tomando únicamente los últimos $k$ dígitos de su parte entera.
		
		\item[k (entero, default=4)] Número de dígitos centrales a extraer. 
		\begin{itemize}
			\item Debe ser un número \textbf{par}.
			\item Se recomienda $k \leq 8$ para evitar errores de precisión de punto flotante en R.
		\end{itemize}
		
		\item[save\_seed (lógico, default=TRUE)] Si es \texttt{TRUE}, almacena el estado final del generador en el entorno global (\texttt{.rng}).
	\end{description}
	
	\section{Detalles del Algoritmo}
	El método genera una secuencia de números iterando sobre los siguientes pasos matemáticos:
	
	\subsection{1. Cuadrado de la semilla}
	Sea $x_i$ la semilla en la iteración $i$, se calcula su cuadrado:
	\begin{equation}
		z = x_i^2
	\end{equation}
	Por ejemplo, para $k=4$ y $x_0=1234$, $z = 1522756$.
	
	\subsection{2. Extracción de dígitos centrales}
	A diferencia de implementaciones clásicas que utilizan sustracciones recursivas, esta función emplea aritmética modular y división entera para aislar los $k$ dígitos centrales. 
	
	Dado $z$, la nueva semilla $x_{i+1}$ se obtiene mediante:
	
	\begin{equation}
		x_{i+1} = \left\lfloor \frac{z}{10^{k/2}} \right\rfloor \pmod{10^k}
	\end{equation}
	
	\textbf{Interpretación:}
	\begin{enumerate}
		\item \textbf{División entera ($\lfloor \cdot \rfloor$):} Elimina los $k/2$ dígitos de la derecha (desplazamiento decimal).
		\item \textbf{Módulo ($\pmod{\cdot}$):} Elimina los dígitos excedentes a la izquierda, conservando los últimos $k$.
	\end{enumerate}
	
	\subsection{3. Normalización}
	El número pseudoaleatorio uniforme $u_i$ se obtiene normalizando la nueva semilla al intervalo $[0, 1)$:
	\begin{equation}
		u_i = \frac{x_{i+1}}{10^k}
	\end{equation}
	
	\section{Persistencia de Estado}
	Si el parámetro \texttt{save\_seed} es \texttt{TRUE}, la función crea o sobrescribe un objeto oculto en el entorno global llamado \texttt{.rng}. Este objeto es una lista con la siguiente estructura, útil para dar continuidad a simulaciones:
	
	\begin{lstlisting}
		.rng = list(
		seed = 5227,       # Ultima semilla generada
		type = "vm",       # Identificador del metodo (von Mises/Neumann)
		parameters = list(k = 4)
		)
	\end{lstlisting}
	
	\section{Limitaciones y Advertencias}
	\begin{itemize}
		\item \textbf{Periodo Corto:} El método de cuadrados medios tiende a tener ciclos cortos comparado con generadores congruenciales lineales o Mersenne Twister.
		\item \textbf{Colapso a Cero:} Si en alguna iteración la semilla toma el valor $0$, toda la secuencia subsecuente será $0$.
		\item \textbf{Precisión Numérica:} R utiliza el estándar IEEE 754 de doble precisión. Para $k > 8$, el valor de $z = \text{seed}^2$ puede exceder $2^{53}$, perdiendo precisión en los dígitos enteros y comprometiendo la aleatoriedad.
	\end{itemize}
	
\end{document}